\documentclass[12pt,a4paper]{article}
\usepackage[utf8]{inputenc}
\usepackage[T1]{fontenc}
\usepackage{amsmath,amsfonts,amssymb,amsthm}
\usepackage[french,english]{babel}
\usepackage{hyperref}

\newtheorem{lemma}{Lemma}[section]
\newtheorem{lemme}{Lemme}[section]
\newtheorem{definition}{Definition}[section]
\newtheorem{theorem}{Theorem}[section]

\title{Draft Setup - Lemma 41}
\author{Charles EDOU NZE \\ \small Independent Researcher}
\date{\today}

\begin{document}

\maketitle

\selectlanguage{english}
\begin{abstract}
This draft setup establishes the exact symbolic typing and formal framework for Lemma 41, pivoting towards the arithmetic rigidity of the Sylow divisor condition as established by Eric Li (arXiv:2606.24872v1). A hypothetical asymmetric deviation from the critical line induces a modification of divisor moments, directly contradicting the strict bounds of the multiplicative large sieve.
\end{abstract}

\section{Symbolic Framework and Typing}
Let $s = \sigma + it$ be a complex variable. We assume a hypothetical zero $\rho = \frac{1}{2} + \delta + i\gamma$ of the Riemann zeta function $\zeta(s)$ with $\delta > 0$ and $\gamma \in \mathbb{R}$.

Let $S(x)$ denote the arithmetic function capturing the Sylow divisor condition. According to Li's resolution of Erdős Problem 768, the asymptotic behavior is tightly bounded:
\begin{equation}
\sum_{n \le x} S(n) \sim C \cdot f(x)
\end{equation}
where $C$ is the exact leading constant.

The deviation $\delta > 0$ induces a twisted moment representation in the multiplicative large sieve space $\mathcal{M}$. We type the sequence of divisor moments as $M_k(\delta) \in \mathcal{H}(\mathcal{M})$, where $\mathcal{H}$ is the functor mapping arithmetic sequences to their Dirichlet series spectral footprint.

\section{Statement of Lemma 41}
\begin{lemma}
Let $\rho = \frac{1}{2} + \delta + i\gamma$ be a non-trivial zero of $\zeta(s)$. The assumption $\delta > 0$ induces a variation in the $k$-th moments $M_k(\delta)$ of the Sylow divisor sums such that for large $x$, the asymptotic bound of the multiplicative large sieve is violated by a factor proportional to $x^\delta$. By the strict arithmetic rigidity of Erdős Problem 768 (Li, 2026), this obstruction is insurmountable, unconditionally forcing $\delta = 0$.
\end{lemma}

\selectlanguage{french}
\section{Cadre Symbolique et Typage}
Soit $s = \sigma + it$ une variable complexe. Supposons l'existence d'un zéro $\rho = \frac{1}{2} + \delta + i\gamma$ de la fonction zêta de Riemann $\zeta(s)$ avec $\delta > 0$ et $\gamma \in \mathbb{R}$.

Soit $S(x)$ la fonction arithmétique capturant la condition du diviseur de Sylow. D'après la résolution du problème 768 d'Erdős par Li, le comportement asymptotique est strictement borné :
\begin{equation}
\sum_{n \le x} S(n) \sim C \cdot f(x)
\end{equation}
où $C$ est la constante principale exacte.

La déviation $\delta > 0$ induit une représentation des moments tordus dans l'espace du grand crible multiplicatif $\mathcal{M}$. Nous typons la séquence des moments de diviseurs comme $M_k(\delta) \in \mathcal{H}(\mathcal{M})$, où $\mathcal{H}$ est le foncteur associant aux suites arithmétiques leur empreinte spectrale via séries de Dirichlet.

\section{Énoncé du Lemme 41}
\begin{lemme}
Soit $\rho = \frac{1}{2} + \delta + i\gamma$ un zéro non trivial de $\zeta(s)$. L'hypothèse $\delta > 0$ induit une variation des moments d'ordre $k$, $M_k(\delta)$, des sommes de diviseurs de Sylow telle que pour un grand $x$, la borne asymptotique du grand crible multiplicatif est violée par un facteur proportionnel à $x^\delta$. Par la stricte rigidité arithmétique du problème 768 d'Erdős (Li, 2026), cette obstruction est insurmontable, forçant inconditionnellement $\delta = 0$.
\end{lemme}

\end{document}
